\emph{सिद्धता-पद्धतींमुळे} वाक्यांमधील निष्पन्नतेचे आणि सुसंगतीचे लक्षणचित्रण करण्यासाठी निव्वळ विन्यासात्मक मार्ग मिळतात. \emph{क्रमवर्ती कलन} ही अशी एक सिद्धता-पद्धती आहे. त्यात \emph{निष्पत्ती} म्हणजे क्रमवर्तींचा सांत वृक्ष असतो (क्रमवर्ती $\Gamma \Sequent \Delta$ म्हणजे~$\Sequent$ या चिन्हाने विभक्त केलेल्या वाक्यांच्या दोन सांत क्रमिका). निष्पत्तीमधील सर्वांत वरच्या क्रमवर्ती \emph{आरंभीच्या क्रमवर्ती} असतात. त्यांची रूपे $!A \Sequent !A$, $\lfalse \Sequent \quad$ आणि बंद पद~$t$ साठी $\quad \Sequent \eq[t][t]$ अशी आहेत. निष्पत्ती योग्य असण्यासाठी इतर प्रत्येक क्रमवर्ती अनेक \emph{निगमन नियमांपैकी} एका नियमानुसार योग्य रीतीने समर्थित असावी लागते. आदिम संयोजक आणि संख्यापक यांच्यासाठीचे नियम जोड्यांनी येतात. प्रत्येक जोडीतील एक नियम क्रमवर्तीच्या डाव्या बाजूला, तर दुसरा उजव्या बाजूला लागू होतो. उदाहरणार्थ, $\Gamma \Sequent \Delta, !A \lif !B$ या क्रमवर्तीचे समर्थन $\RightR{\lif}$ नियमानुसार होत असेल, तर त्याआधीची क्रमवर्ती (म्हणजे \emph{आधारविधान}) $!A, \Gamma \Sequent \Delta, !B$ अशी असली पाहिजे. काही नियम क्रमवर्तीतील वाक्यांचा क्रम किंवा संख्या बदलू देतात. उदाहरणार्थ, \RightR{\Exchange} नियम क्रमवर्तीच्या उजव्या बाजूच्या शेजारील दोन वाक्यांची अदलाबदल करू देतो.

$\quad \Sequent !A$ या क्रमवर्तीची !!a{derivation} असेल, तर $!A$ हे \emph{प्रमेय} आहे असे म्हणतो आणि~$\Proves !A$ असे लिहितो. $\Gamma_0 \Sequent !A$ ची !!a{derivation} असेल आणि $\Gamma_0$ मधील प्रत्येक $!B$ हा~$\Gamma$ मध्ये असेल, तर $!A$ हे~$\Gamma$ पासून \emph{सिद्धतायोग्य} आहे असे म्हणतो आणि $\Gamma \Proves !A$ असे लिहितो. $\Gamma_0 \Sequent \quad$ ची !!a{derivation} असेल आणि $\Gamma_0$ मधील प्रत्येक $!B$ हा~$\Gamma$ मध्ये असेल, तर $\Gamma$ \emph{विसंगत} आहे असे म्हणतो, अन्यथा \emph{सुसंगत} आहे असे म्हणतो. या संकल्पना परस्परांशी संबंधित आहेत. उदाहरणार्थ, $\Gamma \Proves !A$ तेव्हा आणि केवळ तेव्हाच जेव्हा $\Gamma \cup \{\lnot !A\}$ विसंगत असतो. त्यांचा संबंधित चिन्हार्थविषयक संकल्पनांशीही संबंध आहे. उदाहरणार्थ, $\Gamma \Proves !A$ असेल, तर $\Gamma \Entails !A$. $\Gamma$ पासून जे काही सिद्ध करता येते ते~$\Gamma$ पासून तार्किकरीत्या निष्पन्न होते, या सिद्धता-पद्धतींच्या गुणधर्माला \emph{निर्दोषता} म्हणतात. \emph{निर्दोषता प्रमेय} निष्पत्तीतील अनुमानांच्या संख्येवर विगमन करून सिद्ध केले जाते. आरंभीच्या क्रमवर्ती वैध असतात. प्रत्येक अनुमानाची आधारविधाने वैध क्रमवर्ती असतील, तर त्या अनुमानाची निष्कर्ष-क्रमवर्तीही वैध असते.
